\section{Dynamic feature promotion and demotion pipeline}\label{sec:lifecycle}

A feature's tier is a business decision that changes over time: new IP starts cloud-only,
and it moves toward the client only when its secrecy is worth less than the cost and latency
of serving it remotely. Moves toward the client are one-way in practice, because shipped
code must be assumed public, so the pipeline makes demotion (back to the cloud) automatic
and gates promotion to the client on a human reviewer.

\subsection{Tiers}
\begin{itemize}
  \item \textbf{\tier{cloud\_only} (incubation).} Code and weights exist only on servers.
        New or high-value features start here.
  \item \textbf{\tier{hybrid}.} The server stays authoritative, but an obfuscated or
        reduced-quality client fallback is allowed for offline use or poor networks.
  \item \textbf{\tier{client\_ok} (commoditised).} The feature ships in the client package,
        protected only by standard obfuscation \cite{collberg1997taxonomy}, and runs locally
        to save server cost and latency.
\end{itemize}

\begin{figure}[htbp]
\centering
% Figure 2: feature tier state machine with the Section 6.3 thresholds.
\begin{tikzpicture}[
    font=\sffamily\footnotesize,
    st/.style={draw=boxgray, rounded corners=3pt, fill=white, align=left, inner sep=6pt,
               text width=2.9cm, minimum height=12mm},
    arr/.style={-{Stealth[length=5pt]}, draw=boxgray, line width=0.6pt},
    lbl/.style={font=\sffamily\scriptsize, text=textgray, align=center}]
  \node[st] (c) {\textbf{cloud\_only}\\\textcolor{textgray}{Code and weights server-only}};
  \node[st, right=27mm of c] (h) {\textbf{hybrid}\\\textcolor{textgray}{Server-led; client fallback}};
  \node[st, right=27mm of h] (k) {\textbf{client\_ok}\\\textcolor{textgray}{Shipped, obfuscated client}};

  \draw[arr] (c) -- node[lbl, above=1pt, text width=2.5cm] {\(V < 0.65\) and \(S > 0.6\)} (h);
  \draw[arr, draw=promote, line width=1.1pt] (h) -- node[lbl, above=1pt, text width=2.6cm]
      {\(V < 0.40\), on device\\\textcolor{black}{\textbf{human approval}}} (k);
  \draw[arr] (h.south) -- ++(0,-6mm) -| node[lbl, below, pos=0.25] {\(V \ge 0.65\): demote automatically} ($(c.south)+(0.6,0)$);
  \draw[arr] (k.south) -- ++(0,-13mm) -| node[lbl, below, pos=0.25] {\(V \ge 0.65\): next versions return to cloud; shipped code stays public} ($(c.south)+(-0.6,0)$);
\end{tikzpicture}

\caption{Feature tier state machine with the thresholds of \cref{sec:transition}
(\(V\) = IP value, \(S\) = serving pressure; every promotion first waits out a 30-day dwell
time). Only the move to \tier{client\_ok} needs a person to approve it, because shipped code
cannot be recalled; every move back toward the cloud is automatic.}
\label{fig:lifecycle}
\end{figure}

\subsection{Scoring}
The pipeline computes two scores per feature from product and telemetry signals. IP value
\(V\) combines novelty \(N\), share of premium revenue \(R\), competitor parity \(C\) (how
closely public competitors match the output) and whether the method is published,
\(P \in \{0,1\}\):
\begin{equation}\label{eq:value}
  V = 0.35\,N + 0.35\,R + 0.20\,(1 - C) + 0.10\,(1 - P).
\end{equation}
Serving pressure \(S\) compares cost per 1,000 invocations \(c\) and 95th-percentile
latency \(\ell\) to budgets \(c_{\max}\) and \(\ell_{\max}\):
\begin{equation}\label{eq:pressure}
  S = \min\!\left(1,\ \tfrac{1}{2}\,\frac{c}{c_{\max}} + \tfrac{1}{2}\,\frac{\ell}{\ell_{\max}}\right).
\end{equation}
The weights are illustrative; an organisation would set them with its product and legal
teams. Clone risk enters through \(C\) and \(P\): once competitors reach parity or the
method is published, secrecy protects little.

\subsection{Transition policy}\label{sec:transition}
The policy uses a hysteresis band (\(V_{\text{hi}} = 0.65\), \(V_{\text{lo}} = 0.40\)) and
a minimum dwell time of 30 days so that noisy signals do not make a feature oscillate
between tiers (\cref{alg:decide}).

\begin{algorithm}[htbp]
\caption{\Kw{Decide}(feature \(f\), signals \(s\))}\label{alg:decide}
\small
\begin{algorithmic}[1]
\State \(V \gets \Kw{IPValue}(s)\); \(S \gets \Kw{ServingPressure}(s)\) \Comment{\cref{eq:value,eq:pressure}}
\If{\(V \ge V_{\text{hi}}\) \textbf{and} \(f.\text{tier} \ne \tier{cloud\_only}\)} \Return \tier{cloud\_only} \Comment{demote now; reversible} \EndIf
\If{\(f.\text{days\_in\_tier} < \text{MIN\_DWELL}\)} \Return \(f.\text{tier}\) \EndIf
\If{\(f.\text{tier} = \tier{cloud\_only}\) \textbf{and} \(V < V_{\text{hi}}\) \textbf{and} \(S > 0.6\)} \Return \tier{hybrid} \EndIf
\If{\(f.\text{tier} = \tier{hybrid}\) \textbf{and} \(V < V_{\text{lo}}\) \textbf{and} \(s.\text{on\_device\_feasible}\)}
  \State \Return \tier{client\_ok} pending \(\Kw{Approve}(f)\) \Comment{irreversible: human gate}
\EndIf
\State \Return \(f.\text{tier}\)
\end{algorithmic}
\end{algorithm}

\code{lifecycle.py} (\cref{app:lifecycle}) runs one simulated year for a feature whose
novelty and revenue share decay while competitor parity rises and the method is published on
day~201. The feature moves from \tier{cloud\_only} to \tier{hybrid} on day~97
(\(V = 0.648\), \(S = 0.908\)) and to \tier{client\_ok} on day~201 (\(V = 0.36\)), the day
publication removes its remaining secrecy value.

\subsection{Automated packaging}
A promotion to \tier{client\_ok} triggers a build step that compiles the feature for client
targets, applies the shell's standard obfuscation profile, and adds it to the next client
release behind a feature flag. Promotion should never ship the server's exact weights or code
if a distilled or reduced version is enough, since the server version may later become
valuable again or share components with newer cloud-only features.
